\subsection{The neutral net and its index}

The intended neutral current net is
\[
 \A=\A_{D_5,1}\otimes\A_{A_3,1}.
\tag{7.1}\label{eq:basenet}
\]
Its central charge is
\[
 c(D_5,1)+c(A_3,1)=5+3=8.
\tag{7.2}\label{eq:c8}
\]
The discriminant groups of $D_5$ and $A_3$ each have order four, so the
product lattice net has
\[
 \mu_\A=|\det D_5|\,|\det A_3|=4\cdot4=16.
\tag{7.3}\label{eq:mu16}
\]
The isotropic diagonal subgroup generated by $\lambda$ has order four.
The exact lattice identities \eqref{eq:lambda-norm}--\eqref{eq:census}
verify integer spin, evenness and unimodularity of the extension lattice.

\begin{theorem}[Identification once L1/L2 hold]\label{thm:identification}
Assume L1 and L2 produce a local irreducible index-four extension
$\A\subset\B$ with the common stress tensor, with the four grades
generated by $\lambda$, and with finite-dimensional $L_0$ eigenspaces.
Then $\B$ is the $(E_8)_1$ conformal net.
\end{theorem}

\begin{proof}
\emph{Complete rationality and holomorphy.}
The affine lattice net \eqref{eq:basenet} is completely rational, split
and strongly additive.  The order-four local simple-current Q-system has
dimension/index four by Theorem~\ref{thm:Qsystems}.  KLM therefore gives
\[
 \mu_\B=\frac{\mu_\A}{[\B:\A]^2}
 =\frac{16}{4^2}=1.
\tag{7.4}\label{eq:muone}
\]
Hence $\B$ is holomorphic.

\emph{Conformal covariance and split property.}
The common stress tensor is the Osborne--Stottmeister limit, and L2
assumes its covariance acts on the charged primary fields.  Thus the
extension is conformal, not merely M\"obius covariant.  L1 supplies the
energy-bounded generating fields.  Their finite-energy spaces reconstruct
a simple strongly local unitary VOA $V_\B$ by the CKLW/Fredenhagen--J\"or{\ss}
correspondence.  The character below is finite for every
$\operatorname{Im}\tau>0$, so \eqref{eq:traceclass} holds and
Theorem~\ref{thm:split} supplies the split property independently.

\emph{The lattice and character.}
The extension lattice
\[
 L=\bigcup_{k=0}^3(D_5\oplus A_3+k\lambda)
\tag{7.5}\label{eq:gluedlattice}
\]
is even by $\lambda\cdot(D_5\oplus A_3)\subset\Z$ and
$\norm{\lambda}^2=2$.  Its determinant is
\[
 \det L=\frac{\det(D_5)\det(A_3)}{4^2}
 =\frac{16}{16}=1.
\tag{7.6}\label{eq:unimodular}
\]
There is a unique positive-definite even unimodular lattice of rank
eight, namely $E_8$.  Therefore $L\cong E_8$.
Its theta series is the unique normalized weight-four modular form,
\[
 \Theta_L=E_4,
\tag{7.7}\label{eq:thetaE4}
\]
and
\[
 \operatorname{ch}_{V_L}
 =\frac{\Theta_L}{\eta^8}
 =\frac{E_4}{\eta^8}.
\tag{7.8}\label{eq:character}
\]
The exact initial coefficients are
$1,248,4124,34752$; the $248$ weight-one states are the $240$ roots in
\eqref{eq:census} plus eight Cartan currents.

\emph{Uniqueness.}
The reconstructed VOA is simple, unitary, rational, $C_2$-cofinite and
holomorphic: rationality and $C_2$-cofiniteness follow here from the
finite simple-current extension of the positive-definite lattice VOA.
It has $c=8$.  Theorem~\ref{thm:DM} gives
$V_\B\cong V_{E_8}$, and CKLW identifies the generated net with
$\A_{V_{E_8}}$, the level-one $E_8$ current net.
\end{proof}

\begin{remark}[Which hypotheses are presently established]
At finite level, the order-four class, integer spin, census, determinant,
character coefficients, $k=1$ current normalization, and measured
$c=8.004666$ are established.  At continuum level, the neutral
fermion-bilinear currents and stress tensor lie in the
Osborne--Stottmeister theorem class.  What is not established is that the
non-bilinear lattice charged fields satisfy L1/L2 and exhaust the local
scaling algebra.  Without that step, the finite $E_4$ character match is
not a proof of conformal covariance or splitness of the collar limit.
\end{remark}

\begin{corollary}[Local factor type]\label{cor:typeIII}
Under Theorem~\ref{thm:identification}, every nontrivial interval algebra
$\B(I)$ is the hyperfinite type-$\mathrm{III}_1$ factor.
\end{corollary}

\begin{proof}
The $E_8$ level-one current net is split, irreducible, strongly additive
and conformal.  The standard type theorem for nontrivial conformal nets
gives type $\mathrm{III}_1$; splitness and the current-net construction
give hyperfiniteness.
\end{proof}

\section{Downstream corollary cascade}

\begin{corollary}[Charged-field and $\Psi_\lambda$ convergence]
If the final gap box is discharged, the field
$\Psi_{\lambda,N}(f)$ converges on the finite-energy core, is relatively
local, has order-four fusion, and generates the missing odd
$64+64=128$ weight-one states.  The conditional statement in
\emph{$\Psi_\lambda$ at the Quasi-Free Boundary} becomes unconditional
for this collar.
\end{corollary}

\begin{proof}
Apply L1, L2 and Corollary~\ref{cor:Z4}; the state count is
\eqref{eq:census}.
\end{proof}

\begin{corollary}[Type-III KMS uniqueness]
If the final gap box is discharged, the hypotheses of the abstract
holomorphic rotational KMS theorem in
\emph{The Holomorphic Rotational KMS Extension Theorem} hold for the
concrete collar limit.  Thus the rotational $\beta$-KMS state of the
extension is the unique sector Gibbs state, and the canonical
$\Z_4$-invariant extension of the subnet state is
$\omega_\beta\circ E$.
\end{corollary}

\begin{proof}
Theorem~\ref{thm:identification} gives a holomorphic conformal net with
$\mu=1$ and Corollary~\ref{cor:typeIII} gives the type-III local factors.
The expectation is Corollary~\ref{cor:Z4}.  KLM leaves only the vacuum
DHR sector; the rotational KMS classification then leaves only its Gibbs
state.
\end{proof}

\begin{corollary}[P1 and Quillen premises: exact conditional scope]
If, in addition, the determinant-line modular response is identified
with an additive transitive four-channel response and the angular modular
normalization is $\beta_{\rm angle}=2\pi$, then
\[
 c_3=\frac1{4\beta_{\rm angle}}=\frac1{8\pi}.
\tag{8.1}\label{eq:P1}
\]
The KMS state has no independent sector-mixing direction, supplying the
rigidity premise used in the Quillen argument.
\end{corollary}

\begin{proof}
The unique invariant state assigns equal weights to four transitively
permuted channels, hence weight $1/4$ to each.  Multiplying the total
response $1/\beta_{\rm angle}$ by $1/4$ proves \eqref{eq:P1}.
KMS uniqueness removes a second state direction.
\end{proof}

\begin{remark}
The master theorem would not by itself identify the determinant-line
response with the channel response.  That is a separate
P1/Quillen operator-identification premise.  Accordingly the master
theorem closes the \emph{net/KMS rigidity premise}, not the full Quillen
variation formula.
\end{remark}

\section{Per-lemma status}

\renewcommand{\arraystretch}{1.15}
\begin{longtable}{@{}p{0.13\textwidth}p{0.24\textwidth}p{0.16\textwidth}p{0.39\textwidth}@{}}
\toprule
Item & Statement & Status & Reason\\
\midrule
\endhead
S0 & Frozen generator arithmetic:
$\norm{\lambda}^2=2$, $h_\lambda=1$, order four, census
$[52,64,60,64]$ & \mP &
Exact rational lattice calculation; rechecked 10/10.\\
S1 & Neutral bilinear currents and stress tensor have a scaling limit &
\mR &
Osborne--Stottmeister Theorems 4.11, 4.16 and 5.1; the frozen neutral
objects are in the bilinear class.\\
S2a & Continuum $e^\lambda$ has a polynomial energy bound &
\mP &
It is a norm-two, weight-one affine current; Lemma~\ref{lem:continuumEB}
and CKLW.\\
UGF & Edge-projected Green domination and local Sobolev convergence &
\mP &
Theorem~\ref{thm:UGF}: exact particle--hole Schur count plus the sharp
sine comparison gives $K_G=\pi^2/4$ for every $N$.\\
L1 & The finite fields have an $N$-uniform energy bound &
\mP &
Theorem~\ref{thm:L1reduction}, using the proved UGF and norm-one
cocycle/charge shifts.\\
TEL-a & Dyadic edge-mode isometries and charged graph defect &
\mP &
Theorem~\ref{thm:TELa}: explicit second-quantized Shannon embeddings,
$p=3$, $\epsilon=1$, and rate $N^{-2}$ for fields and adjoints.\\
TEL-b & Full-CAR two-edge remainder has an all-$N$ bound &
\mP &
v1022/v1025 inputs plus v1026 CF/DG, one-sided CROSS, validated Fourier
core and full aliases prove $\norm{R_N}_{\HS}<2.995906<3$ for all even
$N\ge16$. Norm only; no one-edge algebra identification or historical
$2.11$ theorem is claimed.\\
L2a & Edge-projected strong convergence on the finite-energy core &
\mP &
The dyadic telescoping series converges by Theorem~\ref{thm:TELa}.\\
ALG-rel & Relative locality with the neutral current net &
\mP &
Theorem~\ref{thm:relativeLocality}: integral root pairings plus strong
and graph convergence pass the exact finite commutator to the limit.\\
ALG-out & Outerness of the limiting $\Z_4$ action &
\mP &
Theorem~\ref{thm:outerness}: the connected multiplicity-free Q-system
gives $\A(I)'\cap\B^{\rm gen}(I)=\C$.\\
ALG-exh & Local algebra convergence and no extra UV operators &
\mG &
Reduced exactly to the two Buchholz--Verch estimates
\eqref{eq:EXHsharp}.\\
L3a & Crossed product commutes with the inductive limit &
\mP &
Theorem~\ref{thm:crossedlimit}, proved isometrically for every finite
group.\\
L3b & Expectations and uniform index bounds commute with the limit &
\mP &
Equations \eqref{eq:coeffE}--\eqref{eq:PP} and
Corollary~\ref{cor:stateE}.\\
L4a & Index-four extension has $\mu=1$ &
\mR &
KLM, after L2 supplies a local irreducible conformal Q-system:
$16/4^2=1$.\\
L4b & Holomorphic $c=8$ extension is $(E_8)_1$ &
\mR &
Even-unimodular rank-eight classification, Dong--Mason and CKLW,
after L1/L2 supply the strongly local VOA/net hypotheses.\\
L4c & Limit local algebras are hyperfinite type $\mathrm{III}_1$ &
\mR &
Standard conformal-current-net factor theorem after L4b.\\
KMS & Unique rotational KMS extension &
\mR &
The downstream theorem applies once L4 is available.\\
P1/Q & Equal four-channel weight and Quillen rigidity premise &
\mR &
Net/KMS part follows; determinant-line response identification remains
separate.\\
\bottomrule
\end{longtable}

\section{The remaining analytic gap}

\begin{center}
\setlength{\fboxsep}{10pt}
\fcolorbox{red!65!black}{red!3}{
\begin{minipage}{0.92\textwidth}
\textbf{Remaining microscopic-exhaustion gaps (status 2026-09-05).}

\smallskip
The Green-form and edge-mode graph parts are no longer assumptions:
Theorems~\ref{thm:UGF} and \ref{thm:TELa} prove them with
$K_G=\pi^2/4$ and dyadic rate $N^{-2}$.  Round-3 certificates close the
interval-TV input, C1 and C2a.  The analytic scalar part of C2b is
$c_{\rm potential}=0.399224424139106\ldots<0.400$.  This does not yet
identify the edge-mode system with the sharply localized scaling algebra
of the sixteen-Majorana $M=1$, $N_y=8$ microscopic collar.

\smallskip\noindent
\textbf{(TEL-B) Relaxed norm obligation closed in round 4.}
For every even $N\ge16$, the now-proved estimates are
\[
 \sqrt N\,\norm{(K'_N-K_N)_{\rm offdiag}}_{\HS}<0.065697,
 \qquad
 \sqrt N\,\norm{(\mathrm{Res}_N-K_N)_{\rm diag}}_{\HS}<0.042959.
\]
The two components are Hilbert--Schmidt orthogonal.  Combined with the
certified scalar contribution, these give $c_E<0.466902<0.542702$.
The one-sided CROSS bound is $0.282637/N$, giving Bound-B squared
$<0.714386$. A 512-node Acb certificate with infinite alias control
proves the weighted Fourier core $<1.783260551$; the actual finite-grid
smooth remainder additionally costs $<0.017333$ in aliases and is
$<1.800594$, not the former unsupported $1.7833$. With the exact jump
norm $\sqrt2$, the unrounded assembly gives $\norm{R_N}_{\HS}<2.995906<3$.
The complete all-$N$ proof is the round-4 supplement in
\texttt{tfpt\_research\_contracts.tex}, executed by
\texttt{v1026\_telb\_round4\_closure.py}. Neither the historical sharper
$2.11$ target nor the absolute-CROSS $0.283/N$ mutant is proved.
The dyadic-increment route is retired, not an alternative premise.

\smallskip\noindent
\textbf{(ALG-EXH) Buchholz--Verch upper inclusion.}
Prove the two estimates \eqref{eq:EXHsharp}; equivalently, construct a
support-preserving microscopic realization of the dyadic embeddings and
prove, for every interval $I$,
\[
 \limsup_N\B_N(I)
 \subset
 W^*\!\left(\A(I),
 \overline{\Psi_{k\lambda}(f)}:
 k\in\Z_4,\ \supp f\subset I\right)
 \subset
 \liminf_N\B_N(I).
\]
Equivalently, no bounded local ultraviolet sequence survives outside
the neutral currents and the four charged grades.  Relative locality
and outerness are already Theorems~\ref{thm:relativeLocality} and
\ref{thm:outerness}; they are not part of this residue.  The binding
ALG2 correction above means that full-CAR Vandermonde interpolation does
not prove this inclusion for $\B_N^{\rm gen}$; FE-GEN or an exact
gauge/code-projected microscopic-algebra theorem is independently required.

\smallskip
For the edge-projected theorem, \textup{(ALG-EXH)} remains the sole hypothesis.
For the original microscopic-collar formulation, the CF/DG norm inputs are
now proved by the Round-4 certificate (v1026), for the fixed-width
$M=1$, $N_y=8$ model and every even $N\ge16$.
This does not supply the microscopic-algebra identification required by
FE-GEN or prove \textup{(ALG-EXH)}.
Once \eqref{eq:EXHsharp} holds, L1/L2 are complete and L3 plus the cited
L4 chain finish the theorem.
\end{minipage}}
\end{center}

\paragraph{Round-7 microscopic twist boundary (2026-09-05).}
The repository checker \texttt{v1033\_charged\_disorder.py} constructs
order-four charged corners in the actual QWZ matter Fock space only after
adding a local four-level flux register to the seam Hamiltonian.
An onsite scalar matter gauge transformation alone cannot change its
Wilson holonomy. The eight-channel quarter twist has charge
$q=(1/4)^8$, hence $h=\lVert q\rVert^2/2=1/4$, not the required
$h_\lambda=1$. Every same-phase integer dressing with norm squared two
has odd fermion parity. The $D_5$ half-twist and $A_3$ inverse-quarter
description additionally requires its diagonal-$U(1)$ quotient and
cocycle/GSO identification; it is not automatically the same microscopic
eight-channel collar. The source's forward seam $i^r$ fixes the Bloch
shift to $k=(2\pi n-r\pi/2)/N$ in the stated convention.
Finite sharp-string energy and nonlocal smooth-twist scaling are
diagnostics, not a conformal-weight or support-preserving FE-GEN proof.
Thus the construction supplies actual corners in an enlarged model,
not the missing MMST upper inclusion. Complete definitions, proofs and
typed regressions are in the round-7 supplement of
\texttt{tfpt\_research\_contracts}.

